\chapter{PART I · 时代背景}

# PART I · 时代背景

\hline\newpage
---

技术革命通常被划分为三个时代：蒸汽机带来的工业革命，电力带来的规模化生产，以及计算机带来的信息化革命。每一次革命的共同特征是：它改变了人类劳动的形式，但始终没有改变一个基本前提——机器服务于人。

人工智能不同。

它正在改变这个前提。当机器开始承担人类的认知劳动——分析、判断、决策、创造——的时候，我们正在经历的，不是又一轮效率升级，而是一次对"人何以为人"的根本性重新定义。

这不是危言耸听。这是每一个身处商业世界中的人，都无法回避的生存问题。

本书开篇，要回答三个看似简单、实则根本的问题：

\textbf{第一，AI究竟是什么？} 它与以往所有的技术浪潮有何本质不同？为什么它不是"风口"，而是文明级的转折点？

\textbf{第二，为什么是德鲁克？} 他的思想诞生于20世纪中叶，为什么在AI时代反而获得了新的生命力？他的哪些框架仍然有效，哪些需要被修正？

\textbf{第三，旧的组织逻辑正在哪些地方失效？} 那些曾经在工业时代和知识经济时代被奉为圭臬的管理常识，在AI时代正在以哪些具体的方式暴露其局限？

这三个问题，是后续六个部分讨论的起点。思维、流程、组织、技能、阻力、未来——所有这些议题，都建立在我们对这三个基础问题的回答之上。

本部分的三章，分别对应这三个问题。它们的目的不是给出终极答案，而是建立一套共同的问题框架——让读者在进入后续更具体的讨论之前，先在认知上做好准备。

因为，在开始重构之前，我们需要先理解：为什么重构是必要的。
